\documentclass[12pt]{article}
\usepackage[margin=1in]{geometry}
\usepackage{enumitem}
\usepackage{titlesec}
\usepackage{parskip}
\usepackage{amsmath}
\usepackage{hyperref}
\usepackage{fontenc}

\title{\textbf{Muller and Strom 1999 Claude Guide}\\
\large M\"uller, Wolfgang C., and Kaare Str\o m, eds. \textit{Policy, Office, or Votes? How Political Parties in Western Europe Make Hard Decisions}.\\
Cambridge: Cambridge University Press, 1999.}
\author{}
\date{}

\begin{document}

\maketitle
\tableofcontents
\newpage

%--------------------------------------------------
\section{Central Claims}
%--------------------------------------------------

\begin{itemize}[leftmargin=*, itemsep=4pt]
    \item \textbf{The core trilemma}: political parties are organizations that can pursue three distinct goals---\textbf{policy} (shaping public policy according to the party's programmatic preferences), \textbf{office} (holding cabinet portfolios and the material/institutional benefits of governing), and \textbf{votes} (maximizing electoral support, valued both instrumentally and for organizational survival/legitimacy). These goals frequently conflict, and no party can simply maximize all three simultaneously.
    \item \textbf{Trade-offs are unavoidable}: pursuing office often requires policy compromise to strike coalition deals; preserving policy purity may mean forgoing office (staying in opposition) or accepting electoral costs; chasing votes by moderating positions toward the median voter can alienate core activists and undercut both policy integrity and internal party cohesion.
    \item \textbf{Parties are not unitary actors}: the volume treats parties as internally differentiated organizations---drawing on the distinction between the party in \textit{public office} (elected officials/ministers), the party in \textit{central office} (leadership and bureaucracy), and the party on \textit{the ground} (members and activists)---whose different actors often have different goal priorities (leaders more office/pragmatic, activists more policy-purist), requiring internal bargaining that shapes how the trilemma is resolved in practice.
    \item \textbf{Institutional context conditions the trade-off}: how sharply parties must choose among policy, office, and votes depends on the electoral system (proportional vs.\ majoritarian, permissiveness), the structure of party competition (number of relevant parties, cleavage structure), and government-formation rules (parliamentary vs.\ presidential, investiture requirements)---echoing and extending Cox-style electoral-institutionalist reasoning to intra-party goal conflict.
    \item \textbf{When parties refuse office}: the book pays particular attention to the puzzle of parties that decline available cabinet seats or coalition offers to preserve policy/ideological integrity or electoral reputation---especially niche, protest, or highly programmatic parties for whom office-holding risks diluting the very identity that sustains their electoral support.
\end{itemize}

%--------------------------------------------------
\section{Theoretical Framework and Antecedents}
%--------------------------------------------------

\begin{itemize}[leftmargin=*, itemsep=4pt]
    \item \textbf{Downsian vote-seeking}: Anthony Downs's median-voter model treats parties as purely vote/office-seeking, converging toward the median voter's position; Muller and Strom's framework relaxes this by allowing parties to genuinely value policy outcomes for their own sake, not merely as instruments to win votes/office.
    \item \textbf{Rikerian office-seeking}: William Riker's coalition theory (minimal winning coalitions) treats parties as office-seeking actors bargaining over the spoils of government; the volume incorporates but complicates this by showing office-seeking is frequently constrained or overridden by policy and electoral considerations.
    \item \textbf{Strom's earlier behavioral theory}: Kaare Strom's 1990 article ``A Behavioral Theory of Competitive Political Parties'' first formalized the policy/office/votes trade-off as a systematic analytic framework (rather than assuming any single-minded goal), arguing that a party's optimal strategy depends on institutional context (e.g., legislative influence available to opposition parties, which affects the attractiveness of staying out of government). This volume extends and empirically tests that framework across many European cases.
    \item \textbf{The tripartite party model}: drawing on Katz and Mair's distinction between the party in public office, in central office, and on the ground, the volume shows that goal conflict is not just about a party's external strategic position but about internal principal-agent problems---leaders (party in public office) may prioritize office access and governing pragmatism, while activists (party on the ground) may prioritize policy purity and threaten defection or internal revolt if leaders compromise too far.
    \item \textbf{Coalition agreements as commitment devices}: the volume examines how written coalition agreements, portfolio allocation rules, and monitoring mechanisms (junior ministers, cabinet committees) function as tools coalition partners use to bind each other to negotiated policy compromises, reducing the risk that the office-seeking behavior of government partners undermines each party's policy goals once in power.
\end{itemize}

%--------------------------------------------------
\section{Comparative Evidence and Case Studies}
%--------------------------------------------------

\begin{itemize}[leftmargin=*, itemsep=4pt]
    \item \textbf{Cross-national case chapters}: contributors examine parties across a wide range of Western European party systems (including Scandinavian, Low Countries, German, and Southern European cases), documenting variation in how readily parties enter coalitions, the policy concessions they accept or refuse, and the conditions under which they choose opposition over office.
    \item \textbf{Niche and anti-establishment parties}: greens, new-left parties, and other programmatically distinctive parties are recurring cases of parties weighing office access against the risk of diluting their core identity and alienating the activist base that sustains their organizational existence.
    \item \textbf{Mainstream/catch-all parties}: larger, more electorally flexible parties are shown to more readily prioritize office and vote-seeking behavior, generally possessing more internal slack to absorb policy compromise without existential organizational risk.
    \item \textbf{Government participation as a strategic choice, not a default}: a recurring empirical finding is that eligibility for office (being part of a winning coalition arithmetically) does not automatically translate into office-seeking behavior---parties frequently and rationally decline available office when the electoral or policy costs are judged too high.
\end{itemize}

%--------------------------------------------------
\section{Core Works Cited}
%--------------------------------------------------

\begin{itemize}[leftmargin=*, itemsep=4pt]
    \item \textbf{Anthony Downs}, \textit{An Economic Theory of Democracy} (1957)---the foundational vote/office-seeking spatial model of party competition that the volume both builds upon and relaxes.
    \item \textbf{William Riker}, \textit{The Theory of Political Coalitions} (1962)---minimal winning coalition theory and office-seeking bargaining, the classic account of coalition formation the volume extends by incorporating policy and electoral goals.
    \item \textbf{Kaare Str\o m}, ``A Behavioral Theory of Competitive Political Parties'' (\textit{American Journal of Political Science}, 1990)---the direct theoretical predecessor establishing the policy/office/votes trade-off framework this volume elaborates and tests.
    \item \textbf{Ian Budge and Michael Laver}, \textit{Party Policy and Government Coalitions} (1992) and related work---formal and empirical analysis of party goals and coalition bargaining over policy.
    \item \textbf{Michael Laver and Norman Schofield}, \textit{Multiparty Government: The Politics of Coalition in Europe} (1990)---a key predecessor synthesizing coalition theory and European comparative evidence.
    \item \textbf{Richard Katz and Peter Mair}, work on party organizational models (party in public office / central office / on the ground)---the organizational framework underlying the volume's treatment of parties as internally differentiated, multi-principal actors.
    \item \textbf{Robert Michels}, \textit{Political Parties} (1911)---the classic statement of oligarchic tension within party organizations (the ``iron law of oligarchy''), relevant background for the volume's leader-activist goal conflict.
\end{itemize}

%--------------------------------------------------
\section{Methodological Contributions and Limitations}
%--------------------------------------------------

\begin{itemize}[leftmargin=*, itemsep=4pt]
    \item \textbf{Disaggregating the unitary-actor assumption}: the volume's central methodological contribution is showing that treating parties as unitary rational actors (standard in both Downsian and Rikerian traditions) obscures important variation driven by intra-party principal-agent dynamics between leaders and activists.
    \item \textbf{Middle-range, comparative-case methodology}: rather than a single formal model tested econometrically, the volume relies on structured, comparable case studies across many European party systems---a design well suited to identifying how institutional context conditions a general behavioral trade-off, though less suited to producing precise, falsifiable point predictions.
    \item \textbf{Limits of generalizability}: the empirical base is concentrated in Western Europe's parliamentary, coalition-prone party systems; applying the office/policy/votes framework to two-party, majoritarian, or non-European (especially non-democratic) party systems requires care, since the institutional conditions that make the trade-off sharply visible (frequent coalition bargaining, proportional representation, viable niche parties) are not universal.
    \item \textbf{Enduring influence}: the policy/office/votes framework became a standard analytic vocabulary in the study of coalition politics and party behavior, widely used in subsequent work on cabinet formation, portfolio allocation, and party strategy.
\end{itemize}

%--------------------------------------------------
\section{Connections to the Broader Literature}
%--------------------------------------------------

\begin{itemize}[leftmargin=*, itemsep=4pt]
    \item \textbf{Cox 1997}: Cox's theory of strategic coordination under different electoral rules explains how many viable parties an electoral system sustains; Muller and Strom take that party landscape as a starting point and ask how the parties that exist within it resolve internal goal conflict, particularly around coalition participation.
    \item \textbf{Cox 1987}: Cox's account of MPs ceding autonomy to disciplined parties for the sake of collective electoral reputation is an institutional-historical instance of exactly the kind of vote-goal-driven organizational discipline that Muller and Strom formalize as one leg of the policy/office/votes trilemma.
    \item \textbf{Lipset and Rokkan 1967}: the cleavage structures that generate and ``freeze'' distinct party identities help explain why some parties (those tightly tied to a specific cleavage-based constituency) face especially sharp policy-purity-vs.-office trade-offs, since compromising core positions risks the party's basis of electoral support.
    \item \textbf{Huber and Shipan 2002}: both works examine how political actors design institutional arrangements (coalition agreements vs.\ statutory delegation) to manage commitment and monitoring problems that arise once decision-making authority is delegated---coalition partners to each other, legislatures to bureaucracies.
    \item \textbf{Dahl 1972}: Dahl's conditions for polyarchy (contestation and inclusiveness) presuppose a party system in which parties compete for office and votes; Muller and Strom's framework adds texture to what that competition actually looks like once multiple, potentially conflicting party goals are taken seriously.
    \item \textbf{Powell 2000}: Powell's ``chain of responsiveness'' linking voters, parties, and policy outcomes depends on parties translating electoral mandates into governing behavior; the policy/office/votes trade-off is precisely the mechanism by which that translation can break down or succeed.
    \item \textbf{Svolik 2012}: Svolik's analysis of power-sharing dilemmas among authoritarian elites (balancing policy control against the security risks of sharing office) offers a structurally analogous, cross-regime-type parallel to the democratic party trilemma of policy, office, and votes.
\end{itemize}

\end{document}
